\begin{afterword}
\summaryhead

The post praises George Orwell as a writer who watched totalitarianism spread and fought it with clear writing, knowing that human evil and muddled thinking go together. It quotes ``Politics and the English Language'' on political speech as the defence of the indefensible, carried out by euphemism, and on plain English, which makes one's own stupid remarks obvious. That, the author says, is the heart of \textsc{Overcoming Bias}.

Evil, the author argues, must keep people's revulsion latent, so it sets out to muddle thinking. The examples are the famines under Stalin and Mao and the scene in \textsc{1984} in which O'Brien tortures Winston to make him say that two and two make five. The author then criticizes the economist Tyler Cowen for doubting that overcoming bias is important, and compares that doubt to the political quietism Orwell saw in Stuart Chase. The post ends by saying that biases shield every evil, that the truth has enemies, and that the fight against our own stupidity is a war.

\responsehead

The post rests on one point that holds: euphemism and vague language help people accept cruelty. Orwell's essay, which the post quotes accurately, is the classic statement of it. The rest of the post widens Orwell's claim and turns it on a critic.

It states Orwell's claims more broadly than Orwell did. Orwell wrote that sloppy language ``makes it easier for us to have foolish thoughts''; the post says Orwell knew ``that muddled language is muddled thinking.'' Orwell described political language ``In our time''; the post says that human evil, ``wherever it exists,'' sets out to muddle thinking, and that biases enshroud every evil. It adds no evidence for the wider claims. Its main examples, the famines under Stalin and Mao, are called ``planned'' and the work of ``deliberate human intent.'' Whether they were intended is disputed, and the Chinese famine is usually traced to policy and to inflated harvest reports.

The post's account of Tyler Cowen does not match Cowen's post. Cowen did not say that overcoming bias is unimportant. He questioned a strong preference for it over other ends, calling bias ``one problem but not the only problem.'' The sentence the post quotes is cut before its qualification, and the post answers it with a question about scope insensitivity, which Cowen did not raise. The likeness to Stuart Chase, whom Orwell described as coming near to calling all abstract words meaningless, is asserted, not shown.

The post says openly why it brings in evil: overcoming bias may not look ``exciting enough'' without ``some clear evil to oppose.'' It then gives the truth enemies, pictures its readers in labor camps, and ends with ``this is a war, and we are soldiers.'' Earlier posts in the same sequence used ``arguments are soldiers'' to describe how politics damages thinking. The war here is against ``ourselves,'' but the post has just set a named critic beside Orwell's opponents.

In short: Orwell's point about language and cruelty, quoted accurately and then widened into a law about all evil, turned against a critic whose post does not say what this one attributes to it, and closed with a call to arms that sits uneasily with the author's own warning against treating arguments as soldiers.
\end{afterword}
